\chapter{Security}
\label{chap:security}

\section{Introduction}
Security in mobile application development is a multi-dimensional challenge that extends beyond the application code to include the runtime environment, data storage mechanisms, and network communication channels. When selecting a development framework, architects must evaluate how the chosen technology influences the application's attack surface.

This chapter presents a comparative analysis of the security postures of \textit{React Native}, \textit{Ionic}, \textit{NativeScript}, \textit{Flutter}, \textit{.NET MAUI}, \textit{Tauri}, \textit{Kotlin Multiplatform (KMP)}, and strictly \textit{Native} development. The comparison focuses on four critical vectors: \textbf{source code obfuscation (reverse engineering resilience)}, \textbf{secure storage capabilities}, \textbf{runtime integrity}, and \textbf{architectural vulnerabilities}.

\section{Architectural Security Classification}
To understand the security implications of each framework, we must first classify them by their execution model, as this dictates their inherent vulnerabilities.

\begin{itemize}
    \item \textbf{WebView-based (Ionic, Tauri):} These applications rely on a browser engine (WebView) to render the UI. The primary security risks are Web vulnerabilities such as Cross-Site Scripting (XSS) and context isolation breaches.
    \item \textbf{Bridge-based (React Native, NativeScript):} These frameworks execute JavaScript (JS) code in a separate runtime and communicate with native modules via a bridge. The logic is interpreted or Just-In-Time (JIT) compiled, making the source code more accessible if not properly bundled.
    \item \textbf{Compiled/Native (Flutter, .NET MAUI, KMP, Native):} These frameworks compile source code into machine code (ARM) or intermediate language (IL). They generally offer higher resilience against reverse engineering compared to JS-based alternatives.
\end{itemize}

\section{Framework-Specific Security Analysis}

\subsection{Native Development (iOS \& Android)}
Native development (Swift/Kotlin) serves as the baseline ``Gold Standard'' for security.
\begin{itemize}
    \item \textbf{Code Protection:} Native apps allow for extensive obfuscation using tools like R8 or ProGuard (Android) and compiler optimizations (iOS). While not impossible to reverse engineer, stripping symbols makes it significantly more time-consuming.
    \item \textbf{API Access:} They have immediate, zero-day access to the latest OS security features (e.g. \textit{CryptoKit}, \textit{BiometricPrompt}) without waiting for third-party wrappers.
    \item \textbf{Vulnerability:} The primary risk is human error—implementing security features inconsistently across the two separate codebases.
\end{itemize}

\subsection{React Native}
React Native is widely adopted but introduces specific risks due to its JavaScript nature.
\begin{itemize}
    \item \textbf{Source Code Exposure:} By default, the business logic is bundled into a JavaScript bundle. Although mostly minified, it is easily readable by attackers. Using the \textit{Hermes} engine helps by pre-compiling JS into bytecode, but it is not as robust as native machine code.
    \item \textbf{Bridge Vulnerabilities:} The communication bridge between JS and Native layers can be intercepted if not properly secured, although the move to the New Architecture (JSI) reduces this serialization overhead and risk.
    \item \textbf{Storage:} It lacks built-in secure storage. Developers must rely on community libraries like \texttt{react-native-keychain}\footnote{https://www.npmjs.com/package/react-native-keychain}, introducing supply-chain risks.
\end{itemize}

\subsection{Flutter}
Flutter creates a unique security profile by bypassing standard OS UI rendering.
\begin{itemize}
    \item \textbf{Reverse Engineering:} Flutter compiles Dart code into a binary snapshot. Reverse engineering these snapshots is notoriously difficult compared to Java/Kotlin decompilation, as the structure is proprietary and lacks standard object headers.~\cite{SEC:FT:01}
    \item \textbf{Data Integrity:} Since Flutter draws its own UI pixels, it is less susceptible to certain OS-level UI tampering or ``screen overlay'' attacks that target standard native views. This is also the reason why password managers might not work in Flutter applications.
    \item \textbf{Weaknesses:} When correctly configured, Flutter presents a relatively small attack surface. However, as with any framework, insecure coding practices at the application level—such as improper key storage or unvalidated user input—can still introduce vulnerabilities.
\end{itemize}

\subsection{Ionic}
Ionic applications are essentially web applications running inside a native container (Capacitor).
\begin{itemize}
    \item \textbf{XSS Risk:} The most critical vulnerability is Cross-Site Scripting. If an attacker injects a script into the WebView, they may gain access to the underlying native bridge, compromising the entire device. \cite{SEC:ION:01}
    \item \textbf{Code Visibility:} As a web-based framework, the source code (HTML/JS/CSS) is packaged in the \texttt{assets} folder. It is trivial to extract and read unless aggressive obfuscation is applied. \cite{SEC:ION:02}
\end{itemize}

\subsection{NativeScript}
NativeScript allows direct access to native APIs via JavaScript without standard wrapper plugins.
\begin{itemize}
    \item \textbf{API Surface:} Because NativeScript can access \textit{any} native API via JS reflection, a compromised JS context has total control over the device (e.g. Camera, Filesystem) without needing specific bridge permission definitions in the code.
    \item \textbf{Obfuscation:} Like React Native, the logic is JavaScript. However, NativeScript offers proprietary JIT compilation which can be slightly harder to inspect than a standard webpack bundle. \cite{SEC:NS:01}
\end{itemize}

\subsection{.NET MAUI (formerly Xamarin)}
MAUI leverages the enterprise-grade .NET ecosystem.
\begin{itemize}
    \item \textbf{Compilation:} On iOS, MAUI uses Ahead-of-Time (AOT) compilation, producing native ARM binaries that are highly secure. On Android, it often relies on Just-In-Time (JIT) and Intermediate Language (IL) assemblies, which are easily decompiled (similar to Java) unless tools like \textit{Dotfuscator} are used.
    \item \textbf{Enterprise Features:} It has strong built-in support for cryptography and authentication libraries authenticated by Microsoft.
\end{itemize}

\subsection{Kotlin Multiplatform (KMP)}
KMP is emerging as a secure alternative by sharing logic while keeping the UI native.
\begin{itemize}
    \item \textbf{Shared Logic Security:} The shared Kotlin code compiles to standard JVM bytecode for Android and an LLVM-based native binary for iOS. This ensures that the iOS security posture is identical to a native Swift app.
    \item \textbf{Auditability:} Since KMP integrates directly into native build systems (Gradle/Xcode), standard security analysis tools work seamlessly, unlike with Flutter or React Native which require specialized scanners.
\end{itemize}

\subsection{Tauri}
Tauri (v2.0+) brings Rust's security guarantees to mobile.
\begin{itemize}
    \item \textbf{Memory Safety:} The backend logic is written in Rust, which is memory-safe by default, eliminating entire classes of vulnerabilities like buffer overflows.
    \item \textbf{Isolation:} Tauri enforces a strict isolation pattern where the frontend (WebView) cannot directly access system APIs. It must pass messages to the Rust backend, which validates all requests. This architectural ``air gap'' makes it significantly more secure than Ionic, despite both using WebViews.
\end{itemize}

\section{Comparative Evaluation}

The following table summarizes the security characteristics of the analyzed frameworks. The ``Reverse Engineering Difficulty'' rates the effort required for an attacker to recover the source logic.

\begin{table}[h!]
\centering
\begin{tabular}{|l|p{3.5cm}|p{3.5cm}|c|}
\hline
\textbf{Framework} & \textbf{Core Language \& Runtime} & \textbf{Source Code Vulnerability} & \textbf{Rev. Eng. Difficulty} \\ \hline
\textbf{Native} & Swift/Kotlin (Native) & Low (Compiled, Stripped) & High \\ \hline
\textbf{Flutter} & Dart (C++ Engine) & Low (Binary Snapshot) & High \\ \hline
\textbf{KMP} & Kotlin (JVM/Native) & Low (Android, ProGuard) / Low (iOS) & High \\ \hline
\textbf{.NET MAUI} & C\# (.NET Runtime) & Medium (IL Decompilation) & Medium-High \\ \hline
\textbf{Tauri} & Rust + JS (WebView) & Low (Backend) / High (Frontend) & High (Backend) \\ \hline
\textbf{React Native} & JS/TS (Hermes/JSI) & High (Bundle Extraction) & Low \\ \hline
\textbf{NativeScript} & JS/TS (V8/JSC) & High (Metadata Reflection) & Low \\ \hline
\textbf{Ionic} & JS/TS (WebView) & High (XSS \& Source) & Very Low \\ \hline
\end{tabular}
\caption{Security Comparison of Mobile Frameworks}
\label{tab:security_comparison}
\end{table}

\section{Conclusion}
From a security perspective, \textbf{Native} development and \textbf{Kotlin Multiplatform} remain the superior choices for high-risk applications (e.g. banking), as they integrate seamlessly with OS-level security protections and obfuscation toolchains. 

\textbf{Flutter} offers a surprisingly robust security profile due to its proprietary binary format, which acts as a ``black box'' to many standard reverse-engineering tools. 

Conversely, JavaScript-based frameworks like \textbf{React Native}, \textbf{NativeScript}, and \textbf{Ionic} require significant additional effort—such as third-party obfuscators (e.g. Jscrambler) and strict CSP (Content Security Policy) configurations—to achieve a comparable security level. \textbf{Tauri} represents a promising middle ground, mitigating the risks of the WebView approach by leveraging a hardened Rust backend.
